% エントリ 20(2026-11-14): 単体版を x86-64 の命令で走らせた確かめ。f の下限
\section{[E20, 2026-11-14] 単体版を x86-64 の命令で走らせた確かめ。$f$ の下限}\label{sec:e20}

E19 の所見 2(単体版が x86 の計算機で止まりうる)は、x86 の挙動を模した版で確かめただけで、x86 では走らせていなかった。ここでは、単体版を x86-64 の命令で実際に走らせた(ユーザー指示「ダウンロードして確かめる」、2026-11-14)。

結論を先に書く。E18 の版は、x86-64 では実際に \texttt{IndexError} で止まる。E19 の版は止まらず、Apple Silicon と同じ値と同じ印を返す。5{,}070{,}000 個の入力で、値の差は最大 $1.9\times10^{-14}$、印の違いは 0 である。E19 の照合で保存した出力(288 万セル)も x86-64 でそのまま再現するので、E19 の照合の結果は x86-64 でも成り立つ。確かめの途中で、x86 とは関係のない誤りを一つ見つけた。$f\lesssim10^{-58}\GeV$ という物理的に意味のない入力で、自己相互作用の率が倍精度をあふれ、自由場の値を印 2 で返していた。$f<10^{-40}\GeV$ を入力の検査で止めるようにした。値を計算する所は E19 から変えていない。まだ確かめていないのは、Linux の x86 の実機と、AVX-512 のある CPU である。

計算は \texttt{calc/e20\_x86.py}(出力 \texttt{calc/out/e20.json})。E18 の版は、E19 と同じく git の 7846eae から読み込む。以下で「いまの版」は、E19 の版に $f$ の下限だけを加えたものを指す。

\subsection{方法}

計算機は E19 と同じ Apple M3 の MacBook である。Rosetta 2(x86-64 の命令を変換して実行するしくみ)の上で、x86-64 の Python と numpy・scipy を動かす。Python は Command Line Tools の 3.9.6(arm64 と x86-64 の両方が入っている)。numpy と scipy は PyPI の x86-64 版の wheel で、作業用のフォルダに展開して読む(conda の環境にもシステムにも入れない。ファイル名と sha256 は \texttt{e20.json})。したがって、これは x86-64 の命令での実行であって、x86 の CPU の実機での実行ではない。

x86-64 は二通りで走らせた。numpy は、CPU が持つ SIMD の命令を調べて、関数ごとに別の実装を選ぶからである。そのままでは numpy が見つける命令は SSE4.2 までで、環境変数 \texttt{ROSETTA\_ADVERTISE\_AVX} を立てると AVX2 と FMA3 までになる(表~\ref{tab:e20env})。NaN を整数に直した結果は、Apple Silicon で 0、x86-64 で $-2^{63}$ で、E19 が模した挙動そのものである。

\begin{table}[H]
\centering\small
\begin{tabular}{lccc}
\toprule
 & Apple Silicon(作った側) & x86-64 & x86-64(AVX2)\\
\midrule
Python & 3.13.11 & 3.9.6 & 3.9.6\\
numpy & 2.3.5 & 2.0.2 & 2.0.2\\
scipy & 1.16.3 & 1.13.1 & 1.13.1\\
numpy が見つける SIMD の命令 & NEON(ASIMD) & SSE4.2 まで & AVX2・FMA3 まで\\
NaN を 64 ビットの整数に直した結果 & 0 & $-2^{63}$ & $-2^{63}$\\
Rosetta 2 による変換 & なし & あり & あり\\
\bottomrule
\end{tabular}
\caption{比べた三つの実行の環境。計算機はどれも同じ Apple M3。}
\label{tab:e20env}
\end{table}

入力は次の組である。いまの版の呼び出しは、すべて警告を例外にして行った。NaN を整数に直すと numpy は警告を出すので、そういう所が残っていれば、どの CPU でもそこで止まる。
\begin{itemize}
\item \textbf{E19 の模した確かめと同じ標本}: 同じ乱数の列の 20 万個 × 6 組($M$ 5--60$\Msun$、$\chi_0$ 0--0.99、年齢 $10^5$--$10^{10}\yr$、$\mu$ 0.6--1.0$\times10^{-12}\eV$。$f$ は $10^{10}$--$10^{19}$、$10^{11}$--$10^{13}$、$10^{17}$--$10^{20}\GeV$ の三通りを 2 組ずつ)。
\item \textbf{値の表の全体とその外}: 30 万個($M$ 1--316$\Msun$、$\chi_0$ 0--0.999、年齢 $10^{4.5}$--$10^{10.3}\yr$、$\mu$ $10^{-13.5}$--$10^{-10.5}\eV$、$f$ $10^{8}$--$10^{20.5}\GeV$)。
\item \textbf{E19 の照合の入力}: seed 1901・1902・1903 の 4{,}800 点 × 200 年齢。E19 のときに Apple Silicon で計算して保存した値と印とも比べる。
\item \textbf{端の値}: 2{,}000 個 × 35 組。一つ目と同じ範囲の入力のうち、一つの引数だけを端の値に置き換える。年齢は 0 と $10^{-30}$--$10^{100}\yr$、$\chi_0$ は 0・$10^{-12}$・0.995・0.9951・0.999・$1-10^{-6}$・$1-10^{-12}$・$1-2^{-53}$、$M$ は $10^{-6}$--$10^{12}\Msun$、$\mu$ は $10^{-30}$--$1\eV$、$f$ は $10^{-40}$--$10^{100}\GeV$。
\item \textbf{広い箱}: 30 万個 × 2 組。引数を同時に広い範囲に振る($M$ $10^{-3}$--$10^{6}\Msun$、$\mu$ $10^{-20}$--$10^{-5}\eV$、年齢 $10^{-3}$--$10^{15}\yr$、$\chi_0$ 0--1。$f$ は $10^{-10}$--$10^{30}\GeV$ と $10^{-40}$--$10^{-10}\GeV$)。
\item \textbf{極端な値}: 2{,}000 個 × 11 組。一つの引数を $10^{\pm300}$ か、正の最小の倍精度($5\times10^{-324}$)に置き換える。
\end{itemize}
E18 の版には、はじめの二つ(7 組)だけを渡した。

\subsection{結果}

\paragraph{E18 の版は x86-64 で止まる}
7 組のうち 4 組で、E18 の版は x86-64 で \texttt{IndexError}(添字が $-2^{63}$)になった(表~\ref{tab:e20old})。Apple Silicon では 7 組とも止まらない。E19 が模した版で見たとおりである。窓の $f=10^{11}$--$10^{13}\GeV$ の 2 組は x86-64 でも止まらない。止まる入力は 20 万個に 2--3 個なので(E19)、同じ範囲でも標本によっては当たらない($10^{10}$--$10^{19}\GeV$ の 2 組目)。SIMD の命令の範囲にはよらない。止まらなかった 3 組では、E18 の版も Apple Silicon と同じ印を返す。

\begin{table}[H]
\centering\small
\begin{tabular}{llccc}
\toprule
標本 & $f$ [GeV] & Apple Silicon & x86-64 & x86-64(AVX2)\\
\midrule
1 組目 & $10^{10}$--$10^{19}$ & ○ & 止まる & 止まる\\
1 組目 & $10^{11}$--$10^{13}$ & ○ & ○ & ○\\
1 組目 & $10^{17}$--$10^{20}$ & ○ & 止まる & 止まる\\
2 組目 & $10^{10}$--$10^{19}$ & ○ & ○ & ○\\
2 組目 & $10^{11}$--$10^{13}$ & ○ & ○ & ○\\
2 組目 & $10^{17}$--$10^{20}$ & ○ & 止まる & 止まる\\
値の表の全体とその外 & $10^{8}$--$10^{20.5}$ & ○ & 止まる & 止まる\\
\bottomrule
\end{tabular}
\caption{E18 の版(git の 7846eae)を、同じ入力で呼んだ結果。「○」は最後まで走った、「止まる」は \texttt{IndexError}。1 組目・2 組目は、E19 の模した確かめで E18 の版・E19 の版に渡した標本。}
\label{tab:e20old}
\end{table}

\paragraph{いまの版は止まらず、同じ値と印を返す}
表~\ref{tab:e20new}。E19 の模した確かめの標本、値の表の全体、E19 の照合の入力、端の値、$f\ge10^{-10}\GeV$ の広い箱(合わせて 4{,}750{,}000 個)では、三つの実行のどれも、例外も警告も出さない。値の差は最大 $1.9\times10^{-14}$ で、印はすべて同じである。値がビットまで同じでない入力は 1{,}918 個(0.04\%)で、差は丸めの大きさである。x86-64 の二通り(SSE4.2 までと AVX2 まで)は、5{,}070{,}000 個の入力すべてで、値がビットまで同じで、印も同じである。Apple Silicon との丸めの違いは、SIMD の命令の範囲からは来ていない。

\begin{table}[H]
\centering\small
\begin{tabular}{lrccc}
\toprule
入力の組 & 入力の数 & 値の差の最大 & ビットまで同じでない数 & 印の違い\\
\midrule
E19 の模した確かめと同じ標本(6 組) & 1{,}200{,}000 & $3.3\times10^{-16}$ & 279 & 0\\
値の表の全体とその外 & 300{,}000 & $3.3\times10^{-16}$ & 33 & 0\\
E19 の照合の入力(3 組) & 2{,}880{,}000 & $1.9\times10^{-14}$ & 1{,}600 & 0\\
端の値(35 組) & 70{,}000 & $1.1\times10^{-16}$ & 2 & 0\\
広い箱、$f\ge10^{-10}\GeV$ & 300{,}000 & $1.1\times10^{-16}$ & 4 & 0\\
\midrule
広い箱、$f<10^{-10}\GeV$ & 300{,}000 & $2.2\times10^{-16}$ & 1 & 0\\
極端な値(10 組) & 20{,}000 & $5.6\times10^{-17}$ & 2 & 0\\
\midrule
合計 & 5{,}070{,}000 & $1.9\times10^{-14}$ & 1{,}921 & 0\\
\bottomrule
\end{tabular}
\caption{いまの版の、Apple Silicon と x86-64 の出力の比較。x86-64 の二通りは同じ数になるので、一つにまとめた。上の 5 行は警告を例外にして呼んだ。$f<10^{-10}\GeV$ の広い箱と極端な値の 2 行は、警告を例外にすると止まる組があるので、その組を警告を無視して呼び直した値で比べた(本文)。極端な値のうち $f=10^{-300}\GeV$ の 1 組は、入力の検査で止まるので入っていない。}
\label{tab:e20new}
\end{table}

\paragraph{E19 の照合の出力の再現}
E19 の照合(\texttt{e19\_validate.py eval})で保存した出力は、4{,}800 点 × 200 年齢 × 3 つの種の 2{,}880{,}000 セルである(値は単精度で保存)。x86-64 の二通りとも、印は全セルで保存した印と同じで、値は単精度に丸めると全セルでビットまで同じである。したがって、E19 の照合(約束を越えるセルが 0)は、x86-64 の出力についてもそのまま成り立つ。細かい年齢の照合(2{,}382 年齢)は出力を保存していないので、この比較には入っていない。

\paragraph{自己診断}
\texttt{share/chi\_sr/selftest.py}(5{,}000 個の入力)は、三つの実行のどれでも通る。値の差の最大は 0、印の違いは 0 である。この計算機のほかの環境(Apple Silicon)でも通る: Python 3.12.15 に numpy 2.5.3・scipy 1.18.1 の環境と、numpy 2.5.4・scipy 1.18.1 の環境。確かめた numpy は 2.0.2 から 2.5.4、scipy は 1.13.1 から 1.18.1、Python は 3.9 から 3.13 である。numpy 1 系では走らせていない。

\paragraph{広い箱と極端な値}
$f\ge10^{-10}\GeV$ の広い箱(30 万個)は、三つの実行とも警告なしで走り、値と印は同じである。$f<10^{-10}\GeV$ の広い箱は、警告を例外にすると、三つの実行とも同じ所で止まる。自己相互作用の率 $\propto\alpha^{11}(\Mpl/f)^4$ が倍精度をあふれるためで、あふれるのは 30 万個のうち 3{,}237 個、どれも $\alpha>10^{6}$ の点である。$\alpha>0.5$ の点はもともと印 3 なので、結果には効かない。警告を無視して呼び直すと、三つとも例外なしで走り、値と印は同じで、全点が $\chi_0$ を返す(差の最大 $7\times10^{-16}$。印は 0・1・3 だけで、64\% が印 3、33\% が印 1)。

極端な値の 11 組のうち 7 組(年齢 $5\times10^{-324}$・$10^{-300}$・$10^{300}\yr$、$\chi_0=5\times10^{-324}$・$10^{-300}$、$\mu=10^{-300}\eV$、$f=10^{300}\GeV$)は、警告なしで走る。$M=10^{-300}\Msun$、$M=10^{300}\Msun$、$\mu=10^{300}\eV$ の 3 組は numpy の警告(0 での割り算、あふれ)を出す。警告を無視すれば三つの実行とも同じ値と印を返し、値は $\chi_0$、印は 0($M=10^{-300}\Msun$)か 3 で、正しい。$f=10^{-300}\GeV$ の組は、いまの版では入力の検査で止まる(次の節)。

\paragraph{速さ}
Rosetta 2 の上での呼び出しの時間は、Apple Silicon の 1.3 倍(SSE4.2 まで)と 2.5 倍(AVX2 まで)だった(表~\ref{tab:e20new} の上の 3 行の入力。組ごとの時間は \texttt{e20.json})。AVX2 までの実行の時間は、走らせるたびに大きく変わる。x86 の実機の速さを示す数ではない。

\subsection{見つかった誤り: $f$ が極端に小さいと、あふれて自由場の値を返す}

極端な値の組のうち $f=10^{-300}\GeV$ は、E19 の版では警告を出し、警告を無視すると誤った値を返した。原因は倍精度のあふれである。自己相互作用の率は $(\Mpl/f)^4$ に比例する(準平衡の判定に使う $\gamma_\infty$ と $\gamma_\mathrm{BH}$)。$\Mpl/f>1.16\times10^{77}$、つまり $f<1.05\times10^{-58}\GeV$ で、この 4 乗が倍精度の上限($1.8\times10^{308}$)を越えて無限大になる。準平衡の占有数 $\epsilon_\mathrm{eq}\propto\sqrt{\gamma_{211}\gamma_\infty}/\gamma_\mathrm{BH}$ は無限大どうしの比で NaN になり、「$\epsilon_\mathrm{eq}$ が自由場の飽和より小さい」という判定が偽になって、点は自由場の側(印 2)に回る。正しくは逆で、自己相互作用が強い極限では雲は育たず、スピンは $\chi_0$ のままである。

\begin{center}\small
\begin{tabular}{lcccccc}
\toprule
$f$ [GeV] & $10^{-10}$ & $10^{-40}$ & $10^{-57}$ & $10^{-58}$ & $10^{-100}$ & $10^{-300}$\\
\midrule
E19 の版 & 0.9000(印 1) & 0.9000(印 1) & 0.9000(印 1) & 0.6005(印 2) & 0.6005(印 2) & 0.6005(印 2)\\
いまの版 & 0.9000(印 1) & 0.9000(印 1) & 止める & 止める & 止める & 止める\\
\bottomrule
\end{tabular}
\end{center}
上の表は $30\Msun$、$\chi_0=0.9$、年齢 $10^9\yr$、$\mu=0.8\times10^{-12}\eV$($\alpha=0.18$)の例である(\texttt{e20\_x86.py smallf})。E19 の版は $f\le10^{-58}\GeV$ で自由場の値 0.6005 を印 2 で返す。

$f$ が $\mu$($10^{-21}\GeV$ の程度)より 37 桁以上小さい範囲で、物理的な意味はなく、集団解析の標本がここに来ることはない。それでも、入力の検査を通る値に対して誤った値を黙って返すのはよくないので、$f<10^{-40}\GeV$ を入力の検査(\texttt{\_check})で止めるようにした(\texttt{ValueError}。プロジェクト版 \texttt{chisr.py} と単体版 \texttt{chi\_sr.py} の両方)。下限は、あふれる所から 18 桁の余裕をとった。下限のところで $\alpha\le0.5$ の点の率があふれるには、$\mu>10^{66}\eV$ が要る。値を計算する所は変えていないので、$f\ge10^{-40}\GeV$ の入力の値と印は E19 の版と同じである。\texttt{export\_chisr.py} に確かめ 6 を加えた: $f=10^{-40}$、$10^{-30}$、$10^{-10}\GeV$ の 2 万個ずつで、単体版もプロジェクト版も警告なしで $\chi_0$ を返し(差の最大 $7\times10^{-16}$、印は 0・1・3 だけ)、下限より小さい $f$ はどちらも \texttt{ValueError} になる。書き出しのほかの確かめ(E19)はすべて通り、表 \texttt{chi\_sr\_tables.npz} と \texttt{selftest.npz} は変わらない。

\begin{result}
\begin{itemize}
\item E18 の版は、x86-64 の命令で実際に \texttt{IndexError} になる(7 組のうち 4 組。Apple Silicon では 0 組)。E19 の所見 2 は、模した版だけでなく実行でも確かめられた。
\item いまの版は、x86-64(SSE4.2 までと AVX2 までの二通り)で止まらず、Apple Silicon と同じ印を返す。5{,}070{,}000 個の入力で値の差は最大 $1.9\times10^{-14}$、印の違いは 0。
\item E19 の照合で保存した 2{,}880{,}000 セルの出力を、x86-64 はそのまま再現する(印は同じ、値は単精度でビットまで同じ)。E19 の照合の結果は x86-64 でも成り立つ。
\item \texttt{selftest.py} は、x86-64 の二通りと、Apple Silicon の numpy 2.3--2.5・scipy 1.16--1.18 で通る。
\item $f\lesssim10^{-58}\GeV$ で自由場の値を印 2 で返す誤り(倍精度のあふれ)を見つけ、$f<10^{-40}\GeV$ を入力の検査で止めた。$f\ge10^{-40}\GeV$ の値と印は E19 と同じ。
\end{itemize}
\end{result}

\begin{caveat}
\begin{itemize}
\item Rosetta 2 による x86-64 の命令の実行であって、x86 の CPU の実機ではない。OS も macOS で、Linux ではない。数学関数の実装(libm)が違えば丸めは変わりうる。値の差は丸めの大きさにとどまるはずだが、閾値のきわの印が変わることはありうる(\texttt{selftest.py} は 5{,}000 個のうち 5 個までの印の違いを許す)。
\item AVX-512 のある CPU(Intel の Xeon など)では、numpy はさらに別の実装を選ぶ。Rosetta 2 は AVX-512 を持たないので、確かめていない。
\item x86-64 で使った numpy と scipy は 2.0.2 と 1.13.1 だけである。CPU と版の両方が作った側と違う一組を確かめたことになる。numpy 1 系は確かめていない。
\item 渡した先の計算機で \texttt{python selftest.py} を走らせる手順は、引き続き必要である。
\item $M$ や $\mu$ が $10^{\pm300}$ のような値のときと、$\alpha\gg1$ で $f$ が極端に小さいときは、numpy の警告(あふれなど)が出る。確かめた範囲では、値と印は正しい($\chi_0$ を、印 0 か 3 で返す)。警告を消してはいない。
\item 入力の検査を通る値の全体で、誤った値を返さないことを示したわけではない。端の値(一つの引数だけ)と広い箱(無作為)で、例外が出ないことと、CPU によらないことを見ただけである。値の正しさを基準の運動論と照合したのは、E19 の範囲($M$ 2--200$\Msun$、$\mu$ $10^{-13}$--$10^{-11}\eV$、$f$ $10^{6}$--$10^{22}\GeV$、年齢 $10^5$--$1.4\times10^{10}\yr$)である。
\end{itemize}
\end{caveat}
